\section{The three algorithms}
\label{sec:algorithms}

Table~\ref{tab:complexity} summarizes the three methods.  The rest of this section
describes each in turn.

\begin{table}[t]
  \caption{Complexity of the three fBM path generation methods.  $N$: time steps;
    $M$: paths; $k$: rSVD target rank.  ``Exact'' means the method samples the
    correct Gaussian distribution on the grid, up to floating-point error; for the
    FFT method this rests on the positivity result of §\ref{sec:circulant}.}
  \label{tab:complexity}
  \centering
  \begin{tabular}{lllll}
    \toprule
    Method & Setup & Per path & Memory & Accuracy \\
    \midrule
    Dense Cholesky   & $O(N^3)$        & $O(N^2)$        & $O(N^2)$  & Exact \\
    Circulant + FFT  & $O(N\log N)$    & $O(N\log N)$    & $O(N)$    & Exact \\
    Low-rank rSVD    & $O(N^2 k)$      & $O(Nk)$         & $O(Nk)$   & Approximate \\
    \bottomrule
  \end{tabular}
\end{table}

\subsection{Algorithm 1: Dense Cholesky}
\label{sec:cholesky}

The fBM covariance matrix $C$ is symmetric positive-definite for all $H \in (0,1)$
and $N \geq 1$, which is sufficient for an exact Cholesky factorization $C = LL^\top$.

\paragraph{Algorithm}
\begin{enumerate}
  \item Build $C \in \mathbb{R}^{N\times N}$ analytically from
        equation~\eqref{eq:fbm-cov}: $O(N^2)$ flops.
  \item Factorize $C = LL^\top$ in place via Cholesky decomposition
        \citep{golub2013matrix}: $N^3/3$ flops.  This is a one-time cost amortized
        over all $M$ paths.
  \item For each path: draw $z \sim \mathcal{N}(0, I_N)$ and compute $\nu Lz$, a
        lower-triangular matrix-vector multiply ($N^2$ flops), which gives the fBM
        path directly; then simulate the stock price path and accumulate the payoff.
\end{enumerate}

The method is exact.  The cost is $O(N^2)$ memory for $L$ (7.6~MB at $N=1000$,
which still fits in the M2's 16~MB L2 cache) and $O(MN^2)$ total work, which
dominates at practical $M$.  This method serves as the correctness baseline:
prices from the other two methods are validated against it.  For the comparison
to be fair, the baseline must do only the work it needs.  An earlier version of my
implementation multiplied by a dense copy of $L$, zeros included, at $2N^2$ flops
per path; §\ref{sec:cholesky-impl} describes the fix and its effect.

\subsection{Algorithm 2: Circulant embedding + FFT (Davies-Harte)}
\label{sec:fft}

The fBM covariance is not Toeplitz, but the covariance of the fGn
increments $\delta W_k$ is (§\ref{sec:toeplitz}). A Toeplitz matrix embeds
into a circulant; circulant matrices are diagonalized by the DFT
(§\ref{sec:circulant}).

\paragraph{Algorithm}
\begin{enumerate}
  \item Compute the fGn autocovariance $\gamma(0), \ldots, \gamma(N-1)$ and form the
        symmetric circulant first row $c \in \mathbb{R}^{2N}$.
  \item Forward FFT: $\lambda = \mathrm{FFT}(c)$; the eigenvalues are all real
        (the circulant is symmetric) and positive (§\ref{sec:circulant}).  The
        code checks the sign and throws if any is negative.
        $O(N\log N)$ flops, computed once.
  \item Pre-compute $s_j = \sqrt{\lambda_j / (2N)}$ for all $j$.
  \item For each path: draw $W \in \mathbb{C}^{2N}$ whose real and imaginary parts
        are i.i.d.\ $\mathcal{N}(0,1)$, scale $W_j \leftarrow s_j W_j$, apply inverse FFT, take the real
        part of the first $N$ entries to obtain fGn increments, cumsum to fBM.
        $O(N\log N)$ flops per path.
\end{enumerate}
FFTW plans for the forward and inverse transforms are created once before the Monte
Carlo loop and reused across all $M$ paths \citep{frigo2005design}.

No clipping is needed.  By the argument of §\ref{sec:circulant}, every eigenvalue
is positive for $H \leq \nicefrac{1}{2}$, so the method is exact on the grid at any
$N$ \citep{craigmile2003simulating}.  The smallest eigenvalue is $0.24\%$ of
$\gamma(0)$ at $H = 0.10$, $N = 252$ (§\ref{sec:stability}): small, but positive by
construction.

\subsection{Algorithm 3: Global low-rank rSVD}
\label{sec:rsvd}

Away from the diagonal $t = s$, the fBM kernel is smooth, so off-diagonal blocks of
$C$ are well-approximated at low rank (§\ref{sec:lowrank}). A global rank-$k$
approximation $C \approx U\,\mathrm{diag}(S)\,U^\top$ replaces $C$ with a compact
factor $L_k = U\,\mathrm{diag}(\sqrt{S}) \in \mathbb{R}^{N\times k}$, dropping the
per-path cost from $O(N^2)$ to $O(Nk)$.

\paragraph{Algorithm}
\begin{enumerate}
  \item Build $C \in \mathbb{R}^{N\times N}$ analytically: $O(N^2)$ flops.
  \item Compute a rank-$(k+p)$ randomized SVD of $C$ with $q = 2$ rounds of
        subspace iteration, re-orthonormalizing after every product (Algorithm~4.4
        of \citet{halko2011finding}; see §\ref{sec:poweriter}): $O(N^2 k)$ flops.
        Retain the top-$k$ singular triplets $(U, S, U^\top)$.
  \item Form the approximate Cholesky factor $L_k = U\,\mathrm{diag}(\sqrt{S})
        \in \mathbb{R}^{N\times k}$: $L_k L_k^\top \approx C$.
  \item For each path: draw $z \sim \mathcal{N}(0, I_k)$, compute $\nu L_k z$
        ($O(Nk)$ flops), cumsum to approximate fBM, simulate payoff.
\end{enumerate}

\paragraph{What it sacrifices.}
Exactness.  The sampled paths have covariance $\nu^2 C_k$ rather than $\nu^2 C$,
with the error controlled by rank $k$.  Section~\ref{sec:errorrank} shows that the
relative Frobenius error falls slowly, from $8.5\%$ at $k = 2$ to $1.2\%$ at
$k = 128$, because the spectrum of the full matrix decays only algebraically
(Figure~\ref{fig:structure}(d)).

This is a global rSVD, a single rank-$k$ approximation of the entire
$N\times N$ matrix.  It must allocate rank budget to both the smooth far-off-diagonal
blocks and the singular near-diagonal region, which drives up the required $k$
at small $H$.  A true hierarchical H-matrix keeps near-diagonal blocks dense and
compresses only far-off-diagonal blocks \citep{hackbusch2015hierarchical}.

\subsection{Other algorithms for further exploration}

\paragraph{True hierarchical H-matrix}
A recursive block-tree H-matrix decomposes $C$ into $O(\log N)$ levels of blocks:
near-diagonal blocks at each level are kept dense, while well-separated blocks at the
same level are approximated at rank proportional to their distance from the diagonal.
The resulting structure supports an H-Cholesky factorization with $O(Nk\log N)$
construction cost and $O(Nk)$ per-path cost, which is better scaling than the global rSVD
at large $N$ \citep{hackbusch2015hierarchical, bebendorf2008hierarchical}.
The complexity gain is a $\log N$ factor over Algorithm~3; it requires implementing
a recursive block-tree and the associated H-matrix arithmetic (§\ref{sec:hmatrix-future}).

\paragraph{Hybrid Scheme for rough Bergomi}
The Hybrid Scheme of \citet{bennedsen2017hybrid} achieves high RMSE reduction
over naïve Euler discretization at the same $O(N\log N)$ cost.  It does so by
exploiting the separability of the Volterra kernel $K(t,s) = (t-s)^{H-1/2}$ near
$t=s$ into a singular analytic part and a smooth remainder.  However, this kernel
governs the rough Bergomi model \citep{bayer2016pricing}, not the fBM-RFSV
model studied here: log-volatility in rough Bergomi is $\int_0^t (t-s)^{H-1/2}dW_s$,
a Volterra integral, whereas in RFSV it is simply $\nu W_t^H$.  Adopting the Hybrid
Scheme would require switching model class.

\paragraph{Quasi-Monte Carlo}
Replacing pseudo-random Gaussian draws with low-discrepancy sequences
\citep{niederreiter1992random} reduces MC standard error from $O(M^{-1/2})$ to near
$O(M^{-1})$, an improvement that can be combined with any of the three algorithms.